\documentclass[10pt,twocolumn]{article}

\usepackage[margin=1in]{geometry}
\usepackage{amsmath,amssymb}
\usepackage{booktabs}
\usepackage{multirow}
\usepackage{graphicx}
\usepackage{hyperref}
\usepackage{url}
\usepackage{microtype}
\usepackage{enumitem}
\usepackage{caption}

\hypersetup{colorlinks=true,linkcolor=blue,citecolor=blue,urlcolor=blue}

\newcommand{\codec}{Vortex-Codec}

\title{\textbf{\codec: A Compressive-Transformer-Based Lossless Codec\\for High Energy Physics Float32 Data}}
\author{[Author Names]\\[Institution(s)]\\\texttt{[correspondence@email.com]}}
\date{Submitted March 2026\\arXiv:XXXX.XXXXX [hep-ex, cs.LG]}

\begin{document}
\maketitle

\begin{abstract}
The Large Hadron Collider (LHC) experiments generate petabytes of IEEE-754 float32 detector data annually. Classical general-purpose compressors (gzip, ZLIB, LZMA, and Zstd) treat these files as opaque byte streams and do not explicitly model the structured entropy of float32 physical observables. We present \codec, a byte-level autoregressive lossless codec that learns float32 entropy structure directly from high-energy physics (HEP) data. \codec combines a Compressive Transformer with Type-Decomposed Token Embeddings (TDTEmbedding), using four learned embedding tables, one per IEEE-754 byte position, and a range-coding backend for exact reconstruction. Evaluated on ATLAS FTAG (2.615 GB) and HEPMC (2.089 GB), \codec achieves 2.35x and 6.34x compression respectively, outperforming all tested classical baselines, including LZMA-9, by 28.0\% and 15.9\% in bits-per-byte (BPD). We also define a cross-experiment transfer-learning protocol and a component-level ablation suite to isolate where gains come from and when IEEE-754-aware embeddings matter.
\end{abstract}

\section{Introduction}
The four major LHC experiments (ATLAS, CMS, LHCb, and ALICE) collectively produce on the order of 15 PB of detector data per year \cite{cern_lhc}. As HL-LHC luminosity ramps, yearly storage pressure is expected to increase substantially.

Current production compression pipelines rely primarily on classical methods: LZMA and ZLIB for bulk HDF5 archival, and Zstd/LZ4 in ROOT-oriented workflows. These algorithms are robust and practical, but they operate at byte-stream level without explicit awareness of IEEE-754 structure. For float32 data, byte positions are not statistically equivalent: sign-plus-exponent bytes are often tightly clustered around physical scales, while low-order mantissa bytes are high entropy. A model that treats all byte positions identically leaves compressible structure unused.

We present \codec, a byte-level autoregressive lossless codec that learns this structure end-to-end. The central architectural contribution is \emph{Type-Decomposed Token Embedding} (TDTEmbedding): four independent 256-entry embedding tables, one for each byte position in a float32 word. This is paired with a Compressive Transformer backbone \cite{compressive_transformer} and Learnable Token Eviction (LTE) \cite{infini_lte} for content-adaptive memory compression.

Concurrent independent work, BOA Constrictor \cite{boa_constrictor}, explores a Mamba-based HEP codec. We cite BOA as concurrent and focus this paper on direct comparisons against strong classical baselines.

\subsection{Headline Results}
\begin{itemize}[leftmargin=1.2em,itemsep=2pt]
  \item \textbf{ATLAS FTAG (2.615 GB, float32 binary):} 3.397 BPD and 2.35x compression. Relative to LZMA-9, BPD improves by 28.0\%.
  \item \textbf{HEPMC (2.089 GB, ASCII event records):} 1.262 BPD and 6.34x compression. Relative to LZMA-9, BPD improves by 15.9\%.
\end{itemize}

\subsection{Contributions}
\begin{itemize}[leftmargin=1.2em,itemsep=2pt]
  \item \textbf{TDTEmbedding:} IEEE-754-aware byte-position embeddings integrated into an end-to-end neural entropy model.
  \item \textbf{Strong lossless compression against classical baselines:} \codec outperforms tested gzip, ZLIB, Bz2, LZMA, and Zstd settings on ATLAS FTAG and HEPMC.
  \item \textbf{Cross-experiment transfer protocol:} a controlled setup for ATLAS pretraining and fine-tuning on CMS, LHCb, ALICE, and CMS NanoAOD.
  \item \textbf{Compressive Transformer with LTE:} content-adaptive compressed memory replacing fixed-stride downsampling.
  \item \textbf{Open-source pipeline:} training, compression, decompression, and evaluation scripts with multi-hardware configs.
\end{itemize}

\section{Related Work}
\subsection{Classical Compression at the LHC}
Production systems use LZ4/Zstd and LZMA/ZLIB in different data-path layers. Schema-level reductions (e.g., xAOD/NanoAOD) reduce stored content but do not replace entropy coding for exact byte reconstruction.

\subsection{Lossy Neural HEP Compression}
BALER \cite{baler} demonstrated learned compression for scientific data but is fundamentally lossy and does not guarantee bit-exact recovery.

\subsection{Neural Lossless Compression}
Neural autoregressive models coupled with entropy coding outperform classical compressors in several non-HEP domains \cite{nncp,gopher_scaling}. GPU-oriented entropy-coding work such as DietGPU motivates practical throughput optimization paths \cite{dietgpu}.

\subsection{Compressive Transformers and LTE}
Compressive Transformer \cite{compressive_transformer} extends effective context with compressed memory. LTE-style token selection \cite{infini_lte} provides a content-aware alternative to uniform striding.

\subsection{Concurrent Work}
BOA Constrictor \cite{boa_constrictor} evaluates a Mamba-based HEP codec. Our study differs in both backbone and byte embedding strategy and defers direct matched-data head-to-head neural comparison.

\section{Architecture}
\codec models a byte stream $b_1, b_2, \ldots, b_N$ autoregressively:
\begin{equation}
P(b_1,\ldots,b_N)=\prod_{t=1}^{N}P(b_t\mid b_{<t}).
\end{equation}
Model probabilities are consumed by a range coder for exact lossless coding. Training minimizes BPD, which directly corresponds to expected codelength.

\subsection{Overview}
Pipeline: input bytes $\rightarrow$ TDTEmbedding $\rightarrow$ positional encoding $\rightarrow$ $N$ transformer blocks (local + compressed-memory attention) $\rightarrow$ linear head (256 logits) $\rightarrow$ softmax probabilities $\rightarrow$ range coder.

\subsection{Type-Decomposed Token Embedding (TDTEmbedding)}
For IEEE-754 float32 streams, byte position carries semantic type information. We use four position-specific embedding tables $E_0, E_1, E_2, E_3\in\mathbb{R}^{256\times d}$ and a learnable scale vector $s\in\mathbb{R}^4$:
\begin{equation}
h_t = E_{t\bmod 4}[b_t]\cdot \mathrm{softmax}(s)_{t\bmod 4}.
\end{equation}

\begin{table}[h]
\centering
\small
\caption{Typical entropy profile by float32 byte position in HEP data.}
\label{tab:tdt_positions}
\begin{tabular}{@{}lll@{}}
\toprule
Byte position & IEEE-754 role & Typical entropy \\
\midrule
0 & Mantissa-low & High (near-uniform) \\
1 & Mantissa-mid & High (near-uniform) \\
2 & Mantissa-high & Medium \\
3 & Sign + exponent & Low (clustered) \\
\bottomrule
\end{tabular}
\end{table}

\subsection{Compressive Attention with LTE}
Each block uses causal local attention plus cross-attention to compressed memory. A per-head gate mixes local and memory outputs. LTE scores historical tokens and selects top-$k$ informative entries (with a differentiable straight-through formulation) instead of fixed-stride compression.

\subsection{Range Coder}
The decoder reproduces the exact byte sequence when provided identical model probabilities in identical order. We verify correctness through byte-for-byte round-trip checks.

\subsection{CATWrapper Dynamic Chunking}
During training, chunk sizes are sampled from \{128, 256, 512\} bytes to encourage multi-scale context use. During inference, the full stream is processed chunk-wise while carrying memory/KV state.

\subsection{Configuration}
\begin{table}[h]
\centering
\small
\caption{Core model configuration used in reported runs.}
\label{tab:model_config}
\begin{tabular}{@{}ll@{}}
\toprule
Parameter & Value \\
\midrule
$d_{model}$ & 512 \\
Layers / heads & 8 / 8 \\
$d_{ff}$ & 2048 \\
Window & 512 bytes \\
Compression rate & 4 \\
Normalization & RMSNorm \\
Precision & fp16 or bf16 \\
\bottomrule
\end{tabular}
\end{table}

\section{Datasets}
We report completed experiments on ATLAS FTAG and HEPMC, with four additional HEP datasets planned.

\begin{table*}[t]
\centering
\small
\caption{Evaluation datasets and current status.}
\label{tab:datasets}
\begin{tabular}{@{}llllll@{}}
\toprule
Dataset & Experiment & Format & Size & Content & Status \\
\midrule
ATLAS FTAG mc-flavtag-ttbar-medium & ATLAS & HDF5 $\rightarrow$ float32 binary & 2.615 GB & Jet kinematics, GN2v01 b-tag, track parameters & Complete \\
HEPMC full event records & MC general & ASCII text + floats & 2.089 GB & Event records, 4-vectors, particle IDs & Complete \\
CMS Run 2 Jet Primary (2012) & CMS & ROOT $\rightarrow$ float32 binary & \~1 GB target & Jet 4-vectors + 24 variables/event & Pending \\
LHCb MC B-decay tracks & LHCb & HDF5 $\rightarrow$ float32 binary & \~500 MB target & Momentum, IP $\chi^2$, PID variables & Pending \\
ALICE TPC Run 3 sample & ALICE & ROOT $\rightarrow$ float32 binary & \~1 GB target & Track-derived float32 features & Pending \\
CMS NanoAOD Run 2 jets & CMS & ROOT $\rightarrow$ float32 binary & \~1 GB target & High-level jet analysis variables & Pending \\
\bottomrule
\end{tabular}
\end{table*}

\subsection{New Dataset Recommendations}
\textbf{CMS Run 2 Jet Primary (2012).} This dataset supports direct lossless-vs-lossy comparison against BALER on the same data and tests transfer across detector geometry.

\textbf{ALICE TPC Run 3.} ALICE TPC has very different structure from jet-centric data and is a strong generalization test for float32-aware embeddings.

\subsection{HEPMC Modality Note}
HEPMC is structured ASCII text, not raw float32 binary. Its strong result is informative for practical compression, but float32-specific architectural claims should rely primarily on binary datasets and binary-specific ablations.

\section{Experimental Setup}
\subsection{Baselines}
Baselines include gzip (L1/L6/L9), ZLIB (L1/L6/L9), Bz2 (L1/L9), LZMA (P6/P9), and Zstd (L1/L3/L9/L19).

\subsection{Training Setup}
Optimization uses AdamW, cosine learning-rate schedule with warmup, gradient clipping (max norm 1.0), mixed precision, and early stopping on validation BPD. Runs were trained on MI300X-class hardware for the main reported checkpoints.

\subsection{Metrics}
Primary metric is BPD (lower is better). We also report compression ratio, encode throughput (MB/s), and relative BPD gain vs LZMA-9.

\section{Results}
\subsection{ATLAS FTAG (2.615 GB)}
\begin{table*}[t]
\centering
\small
\caption{ATLAS FTAG compression results.}
\label{tab:atlas_results}
\begin{tabular}{@{}lcccc@{}}
\toprule
Method & BPD $\downarrow$ & Ratio $\uparrow$ & Speed (MB/s) & vs. LZMA-9 BPD \\
\midrule
\codec (ours) & 3.397 & 2.35x & 1.24 & -28.0\% \\
LZMA-9 & 4.722 & 1.69x & 2.01 & baseline \\
LZMA-6 & 4.741 & 1.69x & 3.38 & +0.4\% \\
Zstd-19 & 4.889 & 1.64x & 5.58 & +3.5\% \\
Zstd-9 & 4.873 & 1.64x & 86.65 & +3.2\% \\
Bz2-9 & 4.947 & 1.62x & 16.35 & +4.8\% \\
Gzip-9 / ZLIB-9 & 4.959 & 1.61x & 13.57 / 13.79 & +5.0\% \\
Zstd-3 & 4.981 & 1.61x & 249.15 & +5.5\% \\
Zstd-1 & 4.980 & 1.61x & 506.29 & +5.5\% \\
Bz2-1 & 5.049 & 1.58x & 16.81 & +6.9\% \\
Gzip-6 / ZLIB-6 & 5.017 & 1.59x & 37.71 / 39.69 & +6.2\% \\
Gzip-1 / ZLIB-1 & 5.268 & 1.52x & 47.55 / 50.92 & +11.6\% \\
\bottomrule
\end{tabular}
\end{table*}

\subsection{HEPMC (2.089 GB)}
\begin{table*}[t]
\centering
\small
\caption{HEPMC compression results.}
\label{tab:hepmc_results}
\begin{tabular}{@{}lcccc@{}}
\toprule
Method & BPD $\downarrow$ & Ratio $\uparrow$ & Speed (MB/s) & vs. LZMA-9 BPD \\
\midrule
\codec (ours) & 1.262 & 6.34x & 2.14 & -15.9\% \\
LZMA-9 & 1.501 & 5.33x & 1.43 & baseline \\
LZMA-6 & 1.525 & 5.25x & 2.09 & +1.6\% \\
Zstd-19 & 1.578 & 5.07x & 2.07 & +5.1\% \\
Bz2-9 & 1.557 & 5.14x & 14.66 & +3.7\% \\
Bz2-1 & 1.655 & 4.83x & 18.23 & +10.3\% \\
Zstd-9 & 1.860 & 4.30x & 56.31 & +23.9\% \\
Zstd-1 & 1.951 & 4.10x & 457.40 & +30.0\% \\
Zstd-3 & 2.087 & 3.83x & 230.26 & +39.0\% \\
Gzip-9 / ZLIB-9 & 2.000 & 4.00x & 9.53 / 9.55 & +33.2\% \\
Gzip-6 / ZLIB-6 & 2.052 & 3.90x & 30.47 / 31.05 & +36.7\% \\
Gzip-1 / ZLIB-1 & 2.461 & 3.25x & 105.08 / 114.39 & +63.9\% \\
\bottomrule
\end{tabular}
\end{table*}

\subsection{Summary}
\begin{table}[h]
\centering
\small
\caption{Current and planned summary across datasets.}
\label{tab:summary}
\begin{tabular}{@{}lcccc@{}}
\toprule
Dataset & Size & BPD & Ratio & vs. LZMA-9 BPD \\
\midrule
ATLAS FTAG & 2.615 GB & 3.397 & 2.35x & -28.0\% \\
HEPMC & 2.089 GB & 1.262 & 6.34x & -15.9\% \\
CMS Run 2 Jets & \~1 GB & TBD & TBD & TBD \\
LHCb MC tracks & \~500 MB & TBD & TBD & TBD \\
ALICE TPC Run 3 & \~1 GB & TBD & TBD & TBD \\
CMS NanoAOD & \~1 GB & TBD & TBD & TBD \\
\bottomrule
\end{tabular}
\end{table}

\section{Ablation Study (Pending)}
The key ablation replaces TDTEmbedding with a single shared embedding table. Running this on both ATLAS (binary float32) and HEPMC (ASCII) is designed to test whether gains are IEEE-754-specific.

\begin{table}[h]
\centering
\small
\caption{Planned ablation matrix.}
\label{tab:ablation}
\begin{tabular}{@{}lll@{}}
\toprule
Variant & ATLAS BPD & HEPMC BPD \\
\midrule
Full \codec & 3.397 (prelim.) & 1.262 (prelim.) \\
w/o TDTEmbedding & TBD & TBD \\
w/o LTE & TBD & TBD \\
w/o Compressive Memory & TBD & TBD \\
w/o CATWrapper & TBD & TBD \\
\bottomrule
\end{tabular}
\end{table}

\section{Cross-Experiment Transfer Learning (Pending)}
We compare two strategies: (A) train from scratch per target dataset; (B) fine-tune from an ATLAS checkpoint. Primary signal is convergence speed to matched BPD.

\begin{table}[h]
\centering
\small
\caption{Transfer-learning experiment design.}
\label{tab:transfer}
\begin{tabular}{@{}lll@{}}
\toprule
Strategy & Target dataset & Primary metric \\
\midrule
Scratch & CMS Run 2 Jets & BPD and ratio vs classical \\
Scratch & LHCb MC tracks & BPD and ratio vs classical \\
Scratch & ALICE TPC Run 3 & BPD and ratio vs classical \\
Fine-tune from ATLAS & CMS Run 2 Jets & Steps to match scratch BPD \\
Fine-tune from ATLAS & LHCb MC tracks & Steps to match scratch BPD \\
Fine-tune from ATLAS & ALICE TPC Run 3 & Steps to match scratch BPD \\
\bottomrule
\end{tabular}
\end{table}

\section{Limitations}
\subsection{Throughput}
On ATLAS, \codec (1.24 MB/s) is slower than LZMA-9 and much slower than fast Zstd settings. The intended deployment target is offline archival, not real-time compression.

\subsection{HEPMC Modality Caveat}
HEPMC is ASCII text, so its gains do not by themselves prove float32-specific architectural benefit.

\subsection{Training Cost}
High-quality checkpoints currently require substantial accelerator resources.

\subsection{Pending Multi-Experiment Results}
CMS, LHCb, ALICE, and NanoAOD experiments, plus full ablations, are still in progress.

\section{Conclusion}
\codec demonstrates that Compressive-Transformer-based neural entropy modeling with IEEE-754-aware embeddings can substantially improve lossless compression on HEP data relative to strong classical baselines. Current completed experiments show 2.35x compression on ATLAS FTAG and 6.34x on HEPMC, with significant BPD improvements over LZMA-9. The pending ablation suite and transfer-learning study are designed to test mechanism-level claims directly: whether byte-position-aware embeddings are specifically beneficial for binary float32 streams and whether pretrained representations transfer across detector modalities.

\bibliographystyle{plain}
\begin{thebibliography}{99}

\bibitem{cern_lhc}
CERN. The Large Hadron Collider.\
\url{https://home.cern/science/accelerators/large-hadron-collider}

\bibitem{compressive_transformer}
J. Rae et al. Compressive Transformers for Long-Range Sequence Modelling. ICLR, 2020. arXiv:1911.05507.

\bibitem{infini_lte}
T. Munkhdalai et al. Leave No Context Behind: Efficient Infinite Context Transformers with Infini-attention / Learnable Token Eviction. arXiv:2404.07143, 2024.

\bibitem{boa_constrictor}
A. Gupta et al. BOA Constrictor: A Mamba-based lossless compressor for High Energy Physics data. arXiv:2511.11337, 2025.

\bibitem{baler}
S. Bengtsson et al. BALER -- Machine Learning Based Compression of Scientific Data. arXiv:2305.02283, 2023.

\bibitem{nncp}
F. Bellard. NNCP: Lossless Data Compression with Transformer.\
\url{https://bellard.org/nncp/}

\bibitem{gopher_scaling}
J. Rae et al. Scaling Language Models: Methods, Analysis \& Insights from Training Gopher. arXiv:2112.11446, 2021.

\bibitem{dietgpu}
Meta Research. DietGPU: GPU-based lossless compression for numerical data.\
\url{https://github.com/facebookresearch/dietgpu}

\end{thebibliography}

\end{document}
